% Please do not change %
\documentclass[11pt]{article} % font size
\usepackage[margin=1in]{geometry} % margins
\usepackage{amsmath,amsthm,amssymb} % for the  common "math symbols"
\usepackage{algorithm}
\usepackage{algpseudocode}
\usepackage[shortlabels]{enumitem} % for enumeration
\usepackage{booktabs} % if you need tables
\usepackage{listings}
\usepackage{color}

\lstset{
  frame=none,
  xleftmargin=2pt,
  stepnumber=1,
  numbers=left,
  numbersep=5pt,
  numberstyle=\ttfamily\tiny\color[gray]{0.3},
  belowcaptionskip=\bigskipamount,
  captionpos=b,
  escapeinside={*'}{'*},
  language=haskell,
  tabsize=2,
  emphstyle={\bf},
  commentstyle=\it,
  stringstyle=\mdseries\rmfamily,
  showspaces=false,
  keywordstyle=\bfseries\rmfamily,
  columns=flexible,
  basicstyle=\small\sffamily,
  showstringspaces=false,
  morecomment=[l]\%,
}
% Feel free to include additional packages (if needed), but make sure it does not override the margin/font requirements.
%
\setlength{\parindent}{0pt} % no indent for new paragraphs.
\setlength{\parskip}{5pt} % distance between paragraphs, which will be used when you have a blank line in your LaTeX code.


%%%%%%%%%%%%%%%%%%%%%%%%%%%%%%%%%%%%%%%%%%%%%%%%%%%%%%%%%%%%%%%%%%%%%%%%%%%%%%
\title{Probability and Computing: \\Formulas}
%
\author{Wira Azmoon Ahmad}
\date{22 Nov 2023}
%
\begin{document}
%
\maketitle

%You may use the usual \LaTeX\ environments to structure your answers:
%\begin{align}
%        \label{basegnn}
%        \mathbf{h}_u^{(t+1)} = \sigma \bigg( \mathbf{W}_\text{self}^{(t)} \mathbf{h}_u^{(t)}  + \mathbf{W}_\text{neigh}^{(t)} \sum_{v \in N(u)} \mathbf{h}_v^{(t)} + \mathbf{b}^{(t)} \bigg). 
%\end{align}


%\begin{table}[h]
%\caption{A simple table.}
%\centering
%\begin{tabular}[t]{lrrrr}
%\toprule
%Graphs & Can & Be  & Fun & Too \\ 
%\midrule 
%$G_1$ & $1$  & $2$ & $3$ & $4$ \\ 
%$G_2$ & $-1$  & $-2$ & $-3$ & $-4$ \\ 
%$G_3$ & $0.1$  & $0.2$ & $0.3$ & $0.4$ \\ 
%\bottomrule
%\end{tabular}
%\label{tab}
%\end{table}

% \section*{Question 1}

\[\forall x, 1+x\leq e^x\]
\[\forall x\in[0,1], 1-\frac{x}{2}\geq e^{-x}\]
\[\forall x\in(0,1),e^{x/2}\leq 1+x\]
\[\forall x\in(0,1),\ln(1+x)\geq \frac{x}{2}\]
\[\forall x\in(-1,1), (1-x)^y \approx e^{-yx}\]
\[\frac{e^x + e^{-x}}{2} \leq e^{x^2/2}\]
\[\left(\frac{n}{k}\right)^k \leq \binom{n}{k} \leq \left(\frac{en}{k}\right)^k\]
Stirling's formula
\[\forall m>0, \sqrt{2\pi m}\left(\frac{m}{e}\right)^m \leq m! 
\leq 2\sqrt{2\pi m}\left(\frac{m}{e}\right)^m\]
Harmonic sum
\[H_n = \sum_{i=1}^n \frac{1}{i} = \ln n + \Theta(1)\]
Basel problem
\[\sum_{i=1}^n \frac{1}{i^2} < \sum_{i = 1}^\infty\frac{1}{i^2} = \frac{\pi^2}{6}
\implies\sum_{i=1}^n \frac{1}{i^2}=\Theta(1)\]
\[\forall x,y, -1<x<y, \frac{x}{1+x}<\frac{y}{1+y}\]

\end{document}
